% Top program acceptance rate function T(theta^T|sigma,OC) (section 4.0.1).
% GENERATED BY generate_figures.py - edit that file, not this one.
% Uniform G, so T is the plain average and the red and green areas are equal.
% xhat^T=0.0488, T=0.54668, curve crosses T at theta=0.2834; areas 0.161694 vs 0.161694
\pgfplotsset{topaxis/.style={
  width=14.8cm, height=7.6cm,
  xmin=-1.75, xmax=2.3, ymin=0.0, ymax=0.9,
  axis x line=bottom, axis y line=left,
  axis line style={-{Stealth[length=7pt]},black},
  clip=false,
  xlabel={$\theta$},
  every axis x label/.style={at={(current axis.right of origin)},anchor=west,xshift=10pt},
  tick label style={font=\small}, label style={font=\small},
  xtick=\empty, tick style={black},
  ytick={0.3000,0.5467},
  yticklabels={$OC$,{$\mathcal{T}(\hat{\theta}^T|\sigma,OC)$}},
}}

\begin{figure}[!ht]
\centering
\begin{tikzpicture}
\begin{axis}[topaxis]
  % shortfall below the average (red) and excess above it (green): equal areas
  \addplot[draw=none,fill=red!30,forget plot] coordinates {
(-1.0000,0.3000) (-0.9671,0.3057) (-0.9342,0.3115) (-0.9013,0.3174) (-0.8684,0.3233)
(-0.8355,0.3292) (-0.8026,0.3352) (-0.7696,0.3412) (-0.7367,0.3472) (-0.7038,0.3533)
(-0.6709,0.3595) (-0.6380,0.3656) (-0.6051,0.3719) (-0.5722,0.3781) (-0.5393,0.3844)
(-0.5064,0.3907) (-0.4735,0.3970) (-0.4406,0.4033) (-0.4077,0.4097) (-0.3748,0.4161)
(-0.3419,0.4226) (-0.3089,0.4290) (-0.2760,0.4355) (-0.2431,0.4420) (-0.2102,0.4485)
(-0.1773,0.4550) (-0.1444,0.4615) (-0.1115,0.4681) (-0.0786,0.4746) (-0.0457,0.4812)
(-0.0128,0.4877) (0.0201,0.4943) (0.0530,0.5008) (0.0859,0.5074) (0.1188,0.5140)
(0.1518,0.5205) (0.1847,0.5271) (0.2176,0.5336) (0.2505,0.5402) (0.2834,0.5467)
(-1.0000,0.5467)} --cycle;
  \addplot[draw=none,fill=green!30,forget plot] coordinates {
(0.2834,0.5467) (0.3171,0.5534) (0.3509,0.5600) (0.3847,0.5667) (0.4184,0.5733)
(0.4522,0.5799) (0.4859,0.5865) (0.5197,0.5931) (0.5535,0.5996) (0.5872,0.6061)
(0.6210,0.6126) (0.6547,0.6190) (0.6885,0.6255) (0.7223,0.6318) (0.7560,0.6382)
(0.7898,0.6445) (0.8235,0.6508) (0.8573,0.6570) (0.8911,0.6632) (0.9248,0.6693)
(0.9586,0.6754) (0.9923,0.6815) (1.0261,0.6875) (1.0598,0.6934) (1.0936,0.6993)
(1.1274,0.7052) (1.1611,0.7110) (1.1949,0.7167) (1.2286,0.7224) (1.2624,0.7280)
(1.2962,0.7336) (1.3299,0.7391) (1.3637,0.7446) (1.3974,0.7499) (1.4312,0.7553)
(1.4650,0.7606) (1.4987,0.7658) (1.5325,0.7709) (1.5662,0.7760) (1.6000,0.7810)
(1.6000,0.5467)} --cycle;
  \draw[gray,dashed] (axis cs:-1.75,0.30000) -- (axis cs:-1.00000,0.30000);
  \draw[black,dashed,thick] (axis cs:-1.75,0.54668) -- (axis cs:1.6,0.54668);
  \draw[gray,dashed] (axis cs:-1.00000,0) -- (axis cs:-1.00000,0.30000);
  \draw[gray,dashed] (axis cs:1.6,0) -- (axis cs:1.6,0.78101);
  \node[anchor=north,font=\small] at (axis cs:-1.00000,-0.012) {$\hat{\theta}^T$};
  \node[anchor=north,font=\small] at (axis cs:1.6,-0.012) {$\bar{\theta}$};
  % winning probabilities: faint where types do not apply to the top program
  \addplot[red!35,thick,smooth] coordinates {
(-1.6000,0.2049) (-1.5640,0.2100) (-1.5281,0.2152) (-1.4921,0.2205) (-1.4562,0.2259)
(-1.4202,0.2313) (-1.3843,0.2368) (-1.3483,0.2424) (-1.3124,0.2481) (-1.2764,0.2538)
(-1.2404,0.2596) (-1.2045,0.2654) (-1.1685,0.2714) (-1.1326,0.2774) (-1.0966,0.2834)
(-1.0607,0.2895) (-1.0247,0.2957)};
  \addplot[red,very thick,smooth] coordinates {
(-1.0000,0.3000) (-0.9559,0.3077) (-0.9119,0.3155) (-0.8678,0.3234) (-0.8237,0.3313)
(-0.7797,0.3394) (-0.7356,0.3475) (-0.6915,0.3556) (-0.6475,0.3639) (-0.6034,0.3722)
(-0.5593,0.3805) (-0.5153,0.3890) (-0.4712,0.3974) (-0.4271,0.4060) (-0.3831,0.4145)
(-0.3390,0.4231) (-0.2949,0.4318) (-0.2508,0.4405) (-0.2068,0.4492) (-0.1627,0.4579)
(-0.1186,0.4666) (-0.0746,0.4754) (-0.0305,0.4842) (0.0136,0.4930) (0.0576,0.5018)
(0.1017,0.5105) (0.1458,0.5193) (0.1898,0.5281) (0.2339,0.5369) (0.2780,0.5456)
(0.3220,0.5543) (0.3661,0.5630) (0.4102,0.5717) (0.4542,0.5803) (0.4983,0.5889)
(0.5424,0.5975) (0.5864,0.6060) (0.6305,0.6144) (0.6746,0.6228) (0.7186,0.6312)
(0.7627,0.6394) (0.8068,0.6477) (0.8508,0.6558) (0.8949,0.6639) (0.9390,0.6719)
(0.9831,0.6798) (1.0271,0.6876) (1.0712,0.6954) (1.1153,0.7031) (1.1593,0.7106)
(1.2034,0.7181) (1.2475,0.7255) (1.2915,0.7328) (1.3356,0.7400) (1.3797,0.7471)
(1.4237,0.7541) (1.4678,0.7610) (1.5119,0.7678) (1.5559,0.7744) (1.6000,0.7810)};
  \node[red,anchor=west,font=\small] at (axis cs:1.7000000000000002,0.7810) {$\theta\mapsto \bar{F}_{\sigma}(\hat{x}^T|\theta)$};
\end{axis}
\end{tikzpicture}
\caption{The top program acceptance rate function. The red curve is the winning probability $\bar{F}_{\sigma}(\hat{x}^T|\theta)$ at the score cutoff \eqref{eq:fixedoc_topdemand}, which meets $OC$ at the marginal top applicant $\hat{\theta}^T$. $\mathcal{T}(\hat{\theta}^T|\sigma,OC)$ is its average over $[\hat{\theta}^T,\bar{\theta}]$: the red and green regions have equal area. $F$ standard normal, $\sigma=2$, $v^T=1$, $G$ uniform on $[-1.6,1.6]$, $OC=0.3$, $\hat{\theta}^T=-1$.}
\label{fig:topacc}
\end{figure}
